
\section{随机变量的数字特征}
\subsection{离散随机变量的期望与方差}
期望的重要性质, 期望与方差的计算


\begin{example}
    已知某类种子每粒发芽的概率都为$ 0.9$, 现播种了$ 1000 $粒, 对于没有发芽的种子,
     每粒需再补种$ 2 $粒, 补种的种子数记为 $X$, 求 $E(X)$, $D(X)$。
\end{example}





\begin{example}
    根据气象预报, 某地区近期有小洪水的概率为$ 0.25$, 
    有大洪水的概率为$ 0.01$.
     该地区某工地上有一台大型设备, 
     遇到大洪水时要损失 $60 000 $元, 遇到小洪水时要损失$ 10 000 $元. 
     为保护设备, 有以下 3 种方案：

\begin{itemize}
    \setlength{\itemsep}{-\itemsep}
    \item 方案 1：运走设备, 搬运费为$ 3 800 $元;
    \item 方案 2：建保护围墙, 建设费为$ 2 000$ 元, 但围墙只能防小洪水; 
    \item 方案 3：不采取措施, 祈祷不发生洪水. 
\end{itemize}
如果你是工地的负责人, 你会选取哪种方案?
\end{example}
\clearpage
\subsection{连续随机变量的期望与方差**}
\clearpage
\subsection{协方差与相关系数**}
\clearpage
\subsection{大数定律与中心极限定理简介**}
\clearpage
